\documentclass[12pt]{article}
\usepackage[utf8]{inputenc}
\usepackage[T1]{fontenc}
\usepackage{hyperref}
\title{Yamben Language Resources for K-12 Curriculum Development}
\author{Compiled by Claude Sonnet 5.5 (claude-sonnet-5-5) for Andr\'as Kornai}
\date{}
\begin{document}
\maketitle

\section{Introduction}
Yamben (ISO 639-3 ynb; Glottocode yamb1257) is a language of Madang Province, Papua New Guinea, spoken in the village of Yambarik, a few hours' hike from Tokain, according to Wikipedia. It was documented by Andrew Pick in the 2010s and classified in 2019 as a probable primary branch of the Madang family, with the exact classification still pending. Joshua Project lists about 90 to 100 speakers, citing Ethnologue. It must not be confused with the separate Yaben language (ybm), which is easy to mix up with it.

\section{Primary Educational Resources}
No primers, textbooks or school programmes were found. The language is not known to be taught in schools.

\section{Linguistic and Reference Materials}
\begin{itemize}
\item Wikipedia, Yamben language: \url{https://en.wikipedia.org/wiki/Yamben_language}
\end{itemize}
The Wikipedia article is the only reference material found. It points to Andrew Pick's documentation work, but I did not locate that work itself in the search.

\section{Religious and Community Resources}
\begin{itemize}
\item Joshua Project, Yamben in Papua New Guinea: \url{https://joshuaproject.net/index.php/people_groups/22573}
\end{itemize}

\section{Digital and Audio Resources}
None found.

\section{Curriculum Development Context}
With a very small speaker community, documentation and community-led materials would have to come before any curriculum. Prospective work should start from Andrew Pick's documentation, which would need to be located through the sources cited in the Wikipedia article.

\end{document}